%% kapa: exam=1 round=28 year=2017 period=2 form=A subject=appraisal_law subject_ko=감정평가관계법규 seq=10 number=10
%% source: Q-Net 감정평가사 기출문제 (cst00309) artlSeq=5199603 file=extracted/1008/5199603/2017년 제28회 감정평가사 1차시험 문제지/2교시 A형.pdf
%% conversion: kordoc 4.24.51 -> tools/gen1.py (auto); review=proofread-sonnet; flags=none
\begin{question}{10}
국토의 계획 및 이용에 관한 법령상 제3종일반주거지역 안에서 건축할 수 있는 건축물을 모두 고른 것은? (단, 조례는 고려하지 않음)
\begin{framedtext}
\noindent\begin{tabular}{@{}p{0.42\linewidth}p{0.42\linewidth}@{}}
ㄱ. 다가구주택 & ㄴ. 아파트 \\
ㄷ. 공중화장실 & ㄹ. 단란주점 \\
ㅁ. 생활숙박시설 & \\
\end{tabular}
\end{framedtext}
\begin{choices}
  \item ㄱ, ㄴ, ㄷ
  \item ㄱ, ㄴ, ㄹ
  \item ㄱ, ㄹ, ㅁ
  \item ㄴ, ㄷ, ㅁ
  \item ㄷ, ㄹ, ㅁ
\end{choices}
\end{question}
\answer{1}
